\documentclass[UTF8,a4paper,11pt,fontset=fandol]{ctexart}
\usepackage[margin=23mm,headheight=15pt]{geometry}
\usepackage{amsmath,amssymb,amsthm,mathtools,bm,mathrsfs}
\usepackage{booktabs,longtable,tabularx,array,enumitem}
\usepackage{xurl}
\usepackage[hidelinks,bookmarksnumbered]{hyperref}
\usepackage{fancyhdr}
\pagestyle{fancy}\fancyhf{}\fancyhead[L]{Ricci夹逼与全局Einstein几何}\fancyfoot[C]{\thepage}
\setlength{\parskip}{4pt}\linespread{1.13}
\setlength{\emergencystretch}{3em}
\setlist{itemsep=2pt,topsep=4pt,parsep=0pt,leftmargin=2em}
\setcounter{tocdepth}{1}
\numberwithin{equation}{section}
\newtheorem{theorem}{定理}[section]
\newtheorem{proposition}[theorem]{命题}
\newtheorem{lemma}[theorem]{引理}
\newtheorem{target}[theorem]{研究目标}
\theoremstyle{definition}
\newtheorem{definition}[theorem]{定义}
\newtheorem{remark}[theorem]{说明}
\newcommand{\Ric}{\operatorname{Ric}}
\newcommand{\Rm}{\operatorname{Rm}}
\newcommand{\Vol}{\operatorname{Vol}}
\newcommand{\diam}{\operatorname{diam}}
\newcommand{\tr}{\operatorname{tr}}
\newcommand{\spec}{\operatorname{Spec}}
\newcommand{\Met}{\operatorname{Met}}
\newcommand{\Id}{\operatorname{Id}}
\newcommand{\op}{\mathrm{op}}
\newcommand{\CP}{\mathbb{CP}}
\newcommand{\R}{\mathbb R}
\newcommand{\Z}{\mathbb Z}
\newcommand{\dd}{\,d\mu_g}
\newcommand{\avgint}{\mathop{\rlap{\,\raisebox{0.1ex}{\tiny$-$}}\int}\nolimits}
\newcommand{\fint}{\mathop{\rlap{\,\raisebox{0.1ex}{\tiny$-$}}\int}\nolimits}
\newcommand{\norm}[1]{\left\lVert#1\right\rVert}
\newcolumntype{L}[1]{>{\raggedright\arraybackslash}p{#1}}

\title{Ricci夹逼下的全局Einstein存在性\\[0.5em]{\large 曲面丛、拓扑障碍与双端几何}}
\author{}
\date{}
\begin{document}
\maketitle
\begin{abstract}
本文研究点态Ricci夹逼与底层光滑流形上Einstein度量存在性的关系。首先由熵、最小Ricci体积和Chern--Gauss--Bonnet公式推导图流形及双曲连接和的定量约束；随后给出固定Dehn-twist曲面丛上的近Einstein度量与全局变分问题。双端连接器部分包含完整Ricci方程、旋转标架的有限区间障碍、正交传递环面作用下的解析障碍，以及带荷稳定中心、局部Einstein芽和全局匹配方程。这些结果区分点态缺陷、积分能量、局部存在与全局存在；一般非对称连接器及曲面丛上的全部Einstein度量仍是开放问题。
\end{abstract}
\begingroup
\small\linespread{1.0}\selectfont\setlength{\parskip}{0pt}
\tableofcontents
\endgroup
\clearpage

\section{问题与归一化}
对每个维数 $n\ge2$，是否存在只依赖 $n$ 的实常数 $c_n<C_n$，使任意闭、连通、光滑实流形 $M^n$ 只要承载一个度量 $g$ 满足
\begin{equation}\label{eq:Q}
 c_ng\le\Ric_g\le C_ng,
\end{equation}
就承载某个光滑 Einstein 度量 $h$，即 $\Ric_h=\lambda h$，$\lambda\in\R$？不要求 $h$ 接近 $g$，不要求 $h$ 与 $g$ 的 Ricci 符号相同；结论针对同一\emph{光滑流形}，不是仅同胚类型。

本文讨论与Yau第26题\cite{Yau1982}相关的上述区间命题，以式\eqref{eq:Q}及其量词作为数学定义。常数不预设正负，输出Einstein度量也不预设符号或与输入的距离。

定义
\begin{equation}\label{eq:defect}
 \delta_\sigma(g)=\norm{g^{-1}\Ric_g-\sigma\Id}_{L^\infty,\op},\qquad
 e_\sigma(M)=\inf_{g\in\Met(M)}\delta_\sigma(g),\qquad \sigma\in\{+1,-1\}.
\end{equation}
$P_n^\sigma$ 表示存在仅依赖 $n$ 的 $\varepsilon_n>0$，使 $\delta_\sigma(g)<\varepsilon_n$ 推出 $M$ 承载某个 Einstein 度量。工作命题满足
\begin{equation}\label{eq:branches}
 Q_n\quad\Longleftrightarrow\quad P_n^+\ \lor\ P_n^-.
\end{equation}
证明只用缩放：$\widehat g=ag$ 时，作为 $(0,2)$ 张量的 Ricci 不变，Ricci 端同态变为 $a^{-1}g^{-1}\Ric_g$。任一非退化区间都有一个非零内点 $\mu$ 及同号小邻域；将 $\Ric\approx\operatorname{sign}(\mu)g$ 的度量按 $1/|\mu|$ 缩放即可进入该邻域。反向直接取以 $\pm1$ 为中心的小区间。

因此，一个固定维数的\emph{证明}只需要完成其中一个符号分支；\emph{完整反驳}要在同一维数击败两个符号分支，可以使用不同流形。每个流形都必须排除任意符号的全部 Einstein 度量。构造 $e_\sigma(M)=0$ 本身是一项独立结论，但它只有与该 $M$ 的全部 Einstein 不存在性相结合，才构成分支反例。

二维可用一致化，三维正分支可用Hamilton的正Ricci结果。以下主要研究四维；其自对偶分解与拓扑恒等式不自动推广到$n\ge5$。

\subsection{曲率范数}
下列 $W^\pm$ 与签名公式在定向四维范畴内使用；非定向情形须另行处理定向双覆盖及等变下降，不能由覆盖的Einstein存在直接推出原流形结论。四维设 $s=\tr_g\Ric$、$E=\Ric-\frac{s}{4}g$。本文件的 $W^\pm$ 使用在 $\Lambda^2_\pm$ 上的曲率\emph{算子}范数；$E$ 和 Ricci 使用张量范数。于是
\begin{align}
8\pi^2\chi(M)&=\int_M\left(|W^+|^2+|W^-|^2+\frac{s^2}{24}-\frac{|E|^2}{2}\right)\dd,\label{eq:CGB}\\
12\pi^2\tau(M)&=\int_M(|W^+|^2-|W^-|^2)\dd,\label{eq:signature}\\
|\Rm_{\rm full}|^2&=4(|W^+|^2+|W^-|^2)+2|E|^2+\frac{s^2}{6}.
\end{align}
若写 $\mathcal R$ 表示完整曲率算子，则 $|\Rm_{\rm full}|^2=4|\mathcal R|^2$。精确 Einstein 时分别有 $\int|\mathcal R|^2=8\pi^2\chi$ 和 $\int|\Rm_{\rm full}|^2=32\pi^2\chi$。Eguchi--Hanson 的 $12\pi^2$ 能量采用前一种规范；后一种为 $48\pi^2$。

$\delta_\sigma\le\delta$ 给出
\begin{equation}
 |s-4\sigma|\le4\delta,\qquad |E|^2\le4\delta^2,\qquad
 |\Ric-\sigma g|^2\le4\delta^2.
\end{equation}
全文不把小 $L^p$ 积分、小加权范数、局部精确解，替换成原题要求的全域无权 Ricci 端同态上确界。

\section{近似输入与全局不存在性}
对一个明确指定的 $M$，反例所需的两部分是
\[
\mathsf P_\sigma(M):\ e_\sigma(M)=0,
\qquad
\mathsf N(M):\ M\text{ 不承载任何Einstein度量}.
\]
分支反例需要同一$M$上的$\mathsf P_\sigma(M)\land\mathsf N(M)$。
正间隙$e_\sigma(M)>0$排除该固定对象的反例性质，却不必给出仅依维数的下界。
另一方面，构造Einstein度量可以判定$\mathsf N(M)$不成立。

\begin{longtable}{L{.22\textwidth}L{.34\textwidth}L{.35\textwidth}}
\toprule 对象 & 已建立的性质 & 开放的全局问题\\\midrule\endhead
指定一般型吹起族 & 精确拓扑间隙与必要退化，见\href{20260916_uniform_gap.tex}{统一间隙专论} & 近似方程的无额外几何假设间隙，或同一无Einstein拓扑上的输入构造\\
固定$A_1$解消 & 满足核心修正和尾项条件时的$e_\sigma(M)=0$构造 & 对其光滑类型判定全部Einstein不存在性，或构造Einstein度量\\
Dehn-twist曲面丛 $E_D$ & 固定光滑对象上的全局近负Einstein输入 & 非适配、非K\"ahler的一般Einstein存在性判断\\
含纯双曲块的extended graph流形 & 在Di Cerbo精确范畴内有全部Einstein不存在性 & 全Ricci缺陷趋零的闭合连接器与全局度量\\\bottomrule
\end{longtable}
曲面丛的$\mathsf P_-$与图流形的$\mathsf N$属于不同对象，不能相加。局部自由$T^2$区域不意味着整个闭流形具有自由环面作用；环面颈与球面连接颈也有不同的拓扑。

\section{扩展图四流形的体积与曲率约束}
\subsection{几何范畴}
采用Di Cerbo文中定义的闭、定向extended graph四流形，满足其双曲块、非纯块、边界仿射粘合及周边子群条件\cite{DiCerbo}。这里的范畴不包含任意有图式分解的流形；无非平凡拼接的闭双曲流形也不属于所讨论的无Einstein对象。

对于含至少一个纯实双曲四维块的相应对象，记各块截面曲率 $-1$ 的体积和为
\[
V_0=\sum_j\Vol_{-1}(P_j)>0.
\]
该文Lemmas 8--9给出 $\tau(M)=0$、$\chi(M)=3V_0/(4\pi^2)$，以及 $\Ric\ge-3g$ 规范下的最小Ricci体积为 $V_0$。这些恒等式与体积下界可进一步约束近Einstein度量。没有纯块的 $\chi=\tau=0$ 情形已受紧夹逼的Chern--Weil障碍限制。

\subsection{近负Einstein序列必须同时饱和体积与Weyl预算}
改用$\Ric\approx-g$归一化，令
\[
V_*=9V_0=12\pi^2\chi(M).
\]
若 $\delta=\delta_-(g)\le1/3$，则
\begin{equation}\label{eq:graphvol}
\boxed{\frac{V_*}{(1+\delta)^2}\le\Vol_gM
\le\frac{V_*}{(1-\delta)^2}.}
\end{equation}
左界来自最小Ricci体积：$\Ric_g\ge-(1+\delta)g$，缩放 $h=(1+\delta)g$ 后 $\Ric_h\ge-h$，其四维体积为 $(1+\delta)^2\Vol_gM$。右界来自CGB与下述点态估计。

对四个Ricci特征值处于 $[-1-\delta,-1+\delta]$，记 $\lambda_i$，则
\[
\frac{s^2}{24}-\frac{|E|^2}{2}
=\frac{(\sum_i\lambda_i)^2}{6}-\frac12\sum_i\lambda_i^2
\ge\frac23(1-\delta)^2,\qquad\delta\le1/3.
\]
可逐变量利用该二次式的凹性把最小化归到区间端点；按高端点数目 $j=0,\dots,4$ 检查五个值即可。结合式\eqref{eq:CGB}和 $8\pi^2\chi=\frac23V_*$，得到式\eqref{eq:graphvol}，并进一步得到
\begin{equation}\label{eq:graphW}
\boxed{0\le\int_M|W|^2\dd\le8\pi^2\chi(M)
\left[1-\left(\frac{1-\delta}{1+\delta}\right)^2\right].}
\end{equation}
所以固定图流形若存在 $\delta_i\to0$，必须有 $\Vol_{g_i}M\to V_*$、$\int|W|^2\to0$；变化拓扑时则是相对 $\chi(M_i)$ 的Weyl能量趋零。另有 $\int|E|^2\le4\delta_i^2\Vol M_i$。

这一约束不产生局部矛盾，因为局部极限可以是真双曲Einstein几何。剩余问题是非双曲粘合拓扑能否集中在很小的Weyl能量区域，同时保持全域Ricci缺陷小。

\subsection{Weyl集中与直径代价}
进一步有不依赖块数的界
\begin{equation}\label{eq:graphconcentration}
 \frac1{\Vol M}\int_M|W|^2\dd\le\frac83\delta,
 \qquad \|W\|_\infty\ge\frac{1-3\delta}{\sqrt3}.
\end{equation}
第一式直接由CGB下界及 $V_*\le(1+\delta)^2\Vol M$ 相减。
第二式来自任意平面的曲率分解
\[
 K_g(\Pi)\le-\frac13+\delta+\frac{|W|}{\sqrt3}:
\]
若最大值界失败，则所有截面严格负，与该候选的
$\mathbb Z^2\subset\pi_1(M)$ 及Preissmann定理矛盾。
所以假想夹逼序列必须把不可消失的Weyl幅度集中到相对体积很小的区域；
例如 $\delta<1/6$ 时，$\{|W|\ge1/(2\sqrt3)\}$ 非空而其体积比例至多 $32\delta$。
这不是自动的能量下界。

若\emph{另有} $\|\nabla W\|_\infty\le L$ 和 $\diam M\le D$，取
$r=\min\{D,1,(4\sqrt3L)^{-1}\}$，$L=0$ 时最后一项视为无穷大。
梯度传播和 $\Ric\ge-7g/6$ 的Bishop--Gromov比较给
\begin{equation}\label{eq:graphdiameter}
 \delta\ge\frac1{128}\frac{v_{-7/18}(r)}{v_{-7/18}(D)},
 \qquad \diam M\ge\sqrt{\frac27}\log(1/\delta)-C_L
 \quad(\delta\text{充分小，}L\text{固定}).
\end{equation}
$v_{-7/18}$为四维常截面曲率模型球体积。具体地，取达到最大值的点，
在半径$r$球上梯度界保证$|W|\ge1/(4\sqrt3)$。由Bishop--Gromov比较，
$\Vol B_r/\Vol M\ge v_{-7/18}(r)/v_{-7/18}(D)$；
与$\fint|W|^2\le8\delta/3$相乘得第一式。
模型体积按$\exp(\sqrt{7/2}D)$增长，给出第二式。
所需$\|\nabla W\|_\infty$界是额外假设，不由零阶Ricci误差推出。

\begin{target}[图流形的二择一结论]
在一个完整核准的上述图流形或明确序列上，构造 $\delta_-(g_i)\to0$；或者建立阻止式\eqref{eq:graphvol}--\eqref{eq:graphW}同时逼近的统一几何代价。固定流形正间隙、固定拓扑族正间隙和仅依维数的正间隙必须分别陈述。
\end{target}
若只证明模型对称类内的正间隙，不能说所有图流形度量已被排除；若真的构造成功，因为外部全部Einstein不存在性已核准，将得到负分支反例，但完整 $Q_4$ 的否定仍需正分支。

\section{双曲连接和的具体间隙：度二BCG与CGB}
这项应用与图流形的相对预算不同，不能因都使用熵而合并为同一定理。
设$(N,g_0)$是闭、定向、实双曲四流形，$\sec(g_0)=-1$，记
$\chi_N=\chi(N)\ge2$、$V_N=\Vol(N,g_0)$。则
\[
 \Ric(g_0)=-3g_0,\qquad V_N=\frac{4\pi^2}{3}\chi_N,
 \qquad\tau(N)=0.
\]
$M=N\#N$经压缩连接颈再折叠两个副本，有度二映射$f:M\to N$。
以$h(g)$表示通用覆盖的体积熵，BCG的度数不等式给
\[
 h(g)^4\Vol(M,g)\ge2\cdot3^4V_N.
\]
若$\Ric_g\ge-3(1+\varepsilon)g$，Bishop--Gromov比较给
$h(g)\le3\sqrt{1+\varepsilon}$，故
\begin{equation}\label{eq:bcg-lower}
 \Vol(M,g)\ge\frac{2V_N}{(1+\varepsilon)^2}.
\end{equation}
这是任意满足Ricci下界的度量的结论，不需要Einstein条件；体积下界来自BCG，
Bishop比较提供的是把熵转成体积所需的上界。

再假定$-3(1+\varepsilon)g\le\Ric_g\le-3(1-\varepsilon)g$，
$0\le\varepsilon\le1/3$。Ricci特征值写成$r_i=-3x_i$，
$x_i\in[1-\varepsilon,1+\varepsilon]$。无迹Ricci仍记为$E$，则
\[
 \frac{s^2}{24}-\frac{|E|^2}{2}
 =\frac{s^2}{6}-\frac{|\Ric|^2}{2}
 =\frac32\left[\left(\sum_i x_i\right)^2-3\sum_i x_i^2\right]
 \ge6(1-\varepsilon)^2.
\]
括号中的二次式对每个变量分别凹，最小值在箱体顶点；检查正端点个数
$0,1,2,3,4$，在上述误差范围最小者为全部$x_i=1-\varepsilon$。
不能沿用早期错误的交叉项公式或在第一项额外加入$1/3$。
由$\chi(M)=2\chi_N-2$与式\eqref{eq:CGB}，得到
\begin{equation}\label{eq:bcg-upper}
 \Vol(M,g)\le
 \frac{2V_N(1-1/\chi_N)}{(1-\varepsilon)^2}.
\end{equation}
两侧相容的必要条件是
\[
 \left(\frac{1-\varepsilon}{1+\varepsilon}\right)^2
 \le1-\frac1{\chi_N}.
\]
因此不存在误差小于下列数值的夹逼输入：
\begin{equation}\label{eq:bcg-threshold}
 \varepsilon_\chi=
 \frac{1-\sqrt{1-1/\chi_N}}{1+\sqrt{1-1/\chi_N}}
 =\frac1{4\chi_N}+O(\chi_N^{-2}).
\end{equation}
$\chi_N\ge2$时$\varepsilon_\chi\le3-2\sqrt2<1/3$，所以推导范围覆盖整个排除区间。
例如$\chi_N=2,4,10$分别给$0.171572875\ldots$、$0.071796770\ldots$、
$0.026334039\ldots$；这是公式的数值展示，不以浮点数代替严格阈值。

\paragraph{归一化与全部Einstein符号。}
在$\Ric\approx-3g$规范，绝对算子缺陷阈值为$3\varepsilon_\chi$；
改为$\widehat g=3g$便有$\delta_-(\widehat g)=\varepsilon$，故
$e_-(N\#N)\ge\varepsilon_\chi$。
同一BCG输入也排除该$M$上的非负Einstein常数：非负Ricci使体积熵为零，
与度二BCG下界矛盾。负Einstein度量可归一化为$\Ric=-3g$，
而式\eqref{eq:bcg-lower}--\eqref{eq:bcg-upper}在$\varepsilon=0$已经矛盾。
这解释了全部Einstein不存在性，而不必假定连接和熵可加。

该间隙依赖$\chi_N$，不是只依赖维数的Yau26间隙。
若让$\chi_N\to\infty$，下界趋零；这既不构造近Einstein度量，
也不证明最佳误差率恰为$1/\chi_N$。固定$N$上的任意紧负夹逼被排除，
变化拓扑上的缺陷趋零仍需真正降低颈部点态误差，不能靠增大外部体积稀释上确界。
阈值计算可用\path{step91_sambusetti_bcg_pinching_threshold.py}重现；
论证只需BCG度数不等式\cite{BCG}和上述点态代数，不依赖连接和熵可加性。

\section{连接区域的相对Chern--Gauss--Bonnet恒等式}
在本节临时使用双曲截面曲率 $-1$、$\Ric=-3g$ 的规范，记物理缺陷 $\mathsf E=\Ric+3g$、$e=\tr\mathsf E$。转回原题取 $h=3g$，则 $\delta_-(h)=(1/3)\norm{g^{-1}\mathsf E}_\infty$，不要混用数值阈值。

设 $N$ 是替换双曲cusp尾部的全部区域，删除尾部总体积为 $V_{\rm rem}$，$V_N=\Vol N$。假设补集保持原精确双曲度量，且拓扑Euler预算对应原块总和。全局CGB减去补集后有
\begin{equation}\label{eq:relative}
\boxed{
2\int_N|W|^2+12(V_N-V_{\rm rem})
=\int_N\left(|\mathsf E^\circ|^2+2e-\frac{e^2}{12}\right).}
\end{equation}
这里必须是匹配后的光滑闭度量；若只对尚有边界的试验域使用，则应保留对应边界转移项。

若外部并非精确双曲，令
$\Psi_g=|W_g|^2+s_g^2/24-|\Ric_g^\circ|^2/2$，
则式\eqref{eq:relative}右侧还要加入
\[
 -2\int_{M\setminus N}(\Psi_g-6)\dd
 +12\{V_{\rm kept}^{\rm ref}-\Vol_g(M\setminus N)\}.
\]
其中 $V_{\rm kept}^{\rm ref}$ 是参考双曲保留区体积。第二项只有在实际与参考外部体积一致时才可删去。
外部曲率与体积域的修正必须同时保留；
$q=e^{2r}q_0$ 仅当 $q_0$ 平坦时才给精确双曲外部，不能把一般 $q_0$ 的衰减残差写成零。

若端同态误差特征值 $x_i\in[-\epsilon,\epsilon]$，二次部分为
\[
\sum_i x_i^2-\frac13\Big(\sum_i x_i\Big)^2,
\]
其绝对值不超过 $4\epsilon^2$，而 $|2e|\le8\epsilon$。所以
\begin{equation}\label{eq:relativebound}
\left|2\int_N|W|^2+12(V_N-V_{\rm rem})\right|
\le(8\epsilon+4\epsilon^2)V_N.
\end{equation}
近零缺陷必满足
\[
V_{\rm rem}-V_N=\frac16\int_N|W|^2+o(V_N).
\]
它是一条相容条件：$V_N$可以小于$V_{\rm rem}$，体积差可以支付Weyl能量。
满足该条件仍需解决水平、形状、混合和荷方程。局部自由$T^2$连接区域的Euler特征
不能代替整个闭流形的Euler特征；任何固定正缺陷也不等价于任意窄区间的反例。

\section{连接器的Ricci方程}
\subsection{Gaussian表示与全部Ricci分量}
令 $g=dr^2+q(r)$，$K=\frac12\partial_rq$、$S=q^{-1}K$、$H=\tr S$。在 $\Ric\approx-3g$ 规范下，残差的全部分量为
\begin{align}
\mathsf E_{rr}&=3-\tr(q^{-1}K_r)+\tr(S^2),\\
\mathsf E_{ri}&=q^{jk}\nabla_jK_{ki}-\partial_iH,\\
\mathsf E_{ij}&=\Ric(q)_{ij}+3q_{ij}-(K_r)_{ij}-HK_{ij}
+2(Kq^{-1}K)_{ij}.\label{eq:fullRic}
\end{align}
等价的法向方程是
\[
H'=3-|S|^2-\mathsf E_{rr}.
\]
混合Codazzi块不能省略。计算范数需通过 $g=LL^{\mathsf T}$ 的Cholesky分解，使用 $L^{-1}\mathsf E L^{-\mathsf T}$ 的谱，而不是坐标分量最大值。

中心Cauchy数据 $(q_0,K_0)$ 应满足标量与动量约束。精确Einstein规范的标量约束为
\[
R_{q_0}=-6+H_0^2-|S_0|^2.
\]
全测地中心 $K_0=0$ 要求 $R_{q_0}=-6$；最小中心 $H_0=0$ 并不要求 $K_0=0$。不能把charged中心擅自替换成全测地中心。

\subsection{正定矩阵表示与端点数据}
为保留全部轨道矩阵和混合方向，令 $B=q_0^{1/2}$，$A=B^{-1}K_0B^{-1}$，$\widehat R=B^{-1}\Ric(q_0)B^{-1}$，写
\[
q=B\exp(X)B.
\]
在 $H_0=0$ 的中心，两侧满足
\[
X(0)=0,\quad X_r(0)=\pm2A,\quad X_{rr}(0)=2(\widehat R+3I).
\]
因 $\frac{d^2}{dr^2}e^X|_{X=0}=X_{rr}+X_r^2$，无须假设矩阵交换，就恢复正确二阶Einstein中心数据。指数表示保证正定，却不控制全Ricci残差，也不证明可延拓到远端。

上述端点一、二阶导数必须由参数化本身满足，而不能在残差计算时单独替换。
周期方向的非线性乘积需要去混叠和分辨率控制，才能把离散残差与连续几何缺陷区分。
完整矩阵与独立分辨率计算分别由
\path{step98_full_matrix_joint_transition.py}、
\path{step98_center_residual_resolution.py}实现。

\subsection{平切片障碍与局部Cauchy问题}
若所有切片平坦，$R_q=0$，标量约束迫使
\[
H^2-|S|^2=6+O(\epsilon),\qquad
H^2\ge9-6\epsilon.
\]
足够小误差时无法穿过 $H=0$。这只排除平切片模型，不排除曲切片、twist或非对角补偿。

另一方面，取非恒定解析周期 $f$，
\[
h_0=d\theta^2+e^{2f(\theta)}dx^2+e^{-2f(\theta)}dy^2,
\qquad R(h_0)=-2(f')^2.
\]
负Yamabe类的归一化代表 $h=\psi^4h_0$ 可满足 $R(h)=-6$，对应
\[
-8\psi''-2(f')^2\psi=-6\psi^5.
\]
它提供Koiso式解析局部Einstein Cauchy数据，说明不能用一个失败的有限参数中心模型排除所有局部中心。它不提供双端cusp匹配、共同长时间延拓或全局闭度量。

\section{长颈、周边格点与荷}
\subsection{半长与误差的共同极限}
在最小中心 $H(0)=0$ 与扩张端 $H(R)=3$ 的受限Gaussian模型中，Raychaudhuri比较给出
\[
1\le k\tanh(kR),\qquad k=\sqrt{1+\epsilon/3}.
\]
若 $\epsilon\to0$，必须 $R\to\infty$；展开相应必要代价为 $\epsilon\gtrsim12e^{-2R}$。这是该端点条件下的法向必要条件，不是全连接器的充分误差率。增加 $R$ 后切向与混合残差可能更坏，不能只展示法向块衰减。

精确外部基准 $q=e^{2r}q_0$ 的残差为切向 $\Ric(q_0)$，物理大小随 $e^{-2r}$ 衰减；最大值仍可能在腰部。$q=\cosh^2(r)q_0$ 也不能保留一般charged中心的Weingarten数据。因此外部redshift不能自动修复内过渡。

\subsection{周边覆盖与尖点提升}
一个有具体拓扑数据的起点是Ratcliffe--Tschantz单尖点24-cell流形的
定向双覆盖 $P=\texttt{24cdc1.1}$\cite[\S\S9,11, Table 2]{RT24}：
\[
 \partial\overline P=T^3,\qquad \chi(P)=2,\qquad
 \Vol_{-1}(P)=\frac{8\pi^2}{3},\qquad
 H_1(P;\Z)=\Z\oplus\Z_3\oplus\Z_{13},\quad H_3(P;\Z)=0.
\]
其最大尖点是原非定向流形体积8的尖点之双覆盖，体积为16。
这里 $\overline P$ 表示截去尖点后的紧核心。
由于边界连通，长正合列和Poincar\'e--Lefschetz对偶给出
\begin{equation}\label{eq:peripheral-rank}
 \dim_{\mathbb Q}\operatorname{im}\bigl(H_1(\partial\overline P;\mathbb Q)
        \longrightarrow H_1(\overline P;\mathbb Q)\bigr)
 =b_1(\overline P)-b_3(\overline P)=1.
\end{equation}
事实上，$H_0(\partial\overline P;\mathbb Q)\to H_0(\overline P;\mathbb Q)$是同构，
故 $H_1(\overline P;\mathbb Q)\to H_1(\overline P,\partial\overline P;\mathbb Q)$满射，
而后者的维数为 $b_3(\overline P)$。因此可选满同态
$\pi_1(P)\to\Z$，在周边群上非零。
这解决了该具体块的周边正秩问题；它还没有指定连接器的仿射粘合、
周边像的整数指数或全部匹配参数，不能当作闭近Einstein度量的构造。

周边格点$(N,1,1)$经cusp深度约$\log N$后得到长度尺度
$(1,N^{-1},N^{-1})$，与rank-two商塌缩中心匹配。
商操作在局部保留度量与Ricci端同态，却把体积与缺陷积分同时缩小；
它不直接减小点态缺陷。

循环覆盖中的cusp数由周边子群在覆盖群中的指数决定。
若覆盖由满同态$\pi_1(P)\to\Z/N\Z$给出，周边群的像为$H_N$，
则该cusp的提升个数是$[\Z/N\Z:H_N]$。
因此整周边在$\Z$中的像为$d\Z$时，$N$重循环覆盖中的提升数为$\gcd(N,d)$；
只有$\gcd(N,d)=1$才得到一个提升cusp。具体双曲块还须满足所用粘合定理的全部假设。

明确地，写尖点度量为 $dt^2+e^{-2t}q_*$，令
$A_*=\Vol(T^3,q_*)$。在深度 $T$ 处，横截体积与尾部体积为
\begin{equation}\label{eq:cusp-tail}
 A(T)=A_*e^{-3T},\qquad
 V_{\rm tail}(T)=\int_T^\infty A_*e^{-3t}\,dt=\frac{A(T)}3.
\end{equation}
在上述单尖点 $N$ 重覆盖中取 $T_N=\log N+T_0$，则
\begin{equation}\label{eq:cusp-cover-scale}
 A_N(T_N)=A_*e^{-3T_0}N^{-2},\qquad
 V_{{\rm rem},N}=\frac{A_*e^{-3T_0}}3N^{-2}.
\end{equation}
两端或多个被替换尖点必须逐个求和，不能把单端尾体积当成总量。
对固定无量纲连接器轮廓，体积、Weyl能量与所有缺陷积分可共同按 $N^{-2}$ 缩放，式\eqref{eq:relative}除以体积后仍不变。真正近零缺陷时 $R$也要增长，匹配深度变为
\begin{equation}\label{eq:RN}
T=\log N-R+O(1),\qquad N\gg e^R.
\end{equation}
所以需要同时控制 $R,N$ 与依 $R$ 变化的轮廓体积和谱常数。把 $N\to\infty$ 当成一切误差自动消失的极限是错误的。

\subsection{荷与Bochner恒等式中的抵消}
对对称矩阵$A$、向量荷$q$和密度$d\mu_\sigma$，设
$Q=\frac12q^{\mathsf T}Aq$、$m_q(X)=q^{\mathsf T}A\,Xq$，则
\begin{equation}\label{eq:charge}
\int m_q(X)d\mu_\sigma=
\int_{\partial\Omega}Q\langle X,\nu\rangle d\mu_{\sigma,\partial}
-\frac12\int q^{\mathsf T}\{X(A)+(\operatorname{div}_\sigma X)A\}q\,d\mu_\sigma.
\end{equation}
它由$XQ=\frac12q^{\mathsf T}X(A)q+m_q(X)$分部积分得到。
连接器中$A=e^{-3\sigma}M^{-1}$，$M$为行列式一的轨道形状矩阵。

带twist的drift曲率产生$+Q|dU|^2$，张力导数却产生
$-Q|dU|^2$抵消；余下的面积与形状梯度相关项没有固定符号。
因此无twist的能量上界不自动适用于变量荷。

\section{曲面丛上的近Einstein构造与全局问题}
\subsection{固定光滑对象的定义与完整慢插值输入}
固定 $F=\Sigma_p$、$B=\Sigma_q$，$p,q\ge2$；
在 $F$ 上取非分离曲线的正Dehn twist $D$，支集外取固定基点。
给定本原满同态 $u:\pi_1(B)\twoheadrightarrow\mathbb Z$，令
\[
 E_D\longrightarrow B,\qquad
 \rho(\gamma)=[D]^{u(\gamma)},\qquad
 \pi_1(E_D)=\pi_1(F)\rtimes_{D_*^{u(\cdot)}}\pi_1(B).
\]
度量参数变化时，$D$的映射类、$p,q,u$和光滑总空间均固定。
这不同于仅在变化拓扑序列上给出小误差。

固定双曲纤维领圈
$h_0=du_0^2+a(u_0)^2ds^2$，$a=a_0\cosh u_0$，
取端部常值 $0,1$ 的光滑函数 $\zeta$，写
\[
 h_z=du_0^2+a^2(ds+z\zeta' du_0)^2,
 \quad h_1=D^*h_0,
 \quad K_{h_z}=-1,
 \quad dA_{h_z}=dA_{h_0}.
\]
度量在领圈边界匹配；非整数 $z$ 对应的局部坐标变换不必是闭纤维的全局同痕。
在基底捏缩一条非分离曲线，长度 $\ell\to0$，置
\[
 L=\tfrac12\operatorname{arsinh}(1/\sinh(\ell/2)),\quad
 b(t)=\ell\cosh t,\quad z(t)=\eta(t/L),\quad
 G_\ell=dt^2+b^2d\theta^2+h_{z(t)}.
\]
$\eta$在两端常值；外部取双曲乘积，通过右端$D$粘接，得到整个固定$E_D$的光滑闭度量。
记$j=\tanh u_0$、$k=\tanh t$、$Q=a\zeta'$、$v=Qz'$。
在正交余标架$(dt,b\,d\theta,du_0,a(ds+z\zeta'du_0))$中，全部缺陷为
\begin{equation}\label{eq:bundleRicci}
 (\Ric_{G_\ell}+G_\ell)_{\rm ON}=
 \begin{pmatrix}
 -v^2/2&0&0&(v_{u_0}+2jv)/2\\
 0&0&0&0\\
 0&0&-v^2/2&-(v_t+kv)/2\\
 (v_{u_0}+2jv)/2&0&-(v_t+kv)/2&v^2/2
 \end{pmatrix}.
\end{equation}
形算子协变导数的二次项包含标架旋转，不能只对矩阵条目普通求导。
令 $A_0=\|Q\|_\infty$、$A_1=\|Q'+2jQ\|_\infty$、
$M_i=\|\eta^{(i)}\|_\infty$，则
\begin{equation}\label{eq:bundlebound}
 \begin{aligned}
 \delta_-(G_\ell)&\le\frac{A_0^2M_1^2}{2L^2}\\
 &\quad+\frac1{2L}\sqrt{A_1^2M_1^2+A_0^2(M_1+M_2/L)^2}\longrightarrow0.
 \end{aligned}
\end{equation}
误差支持之外为精确乘积；该式是全局无权端同态界，不是体积小导致的积分误差。
同时
\[
 \begin{gathered}
 \Vol(G_\ell)=V_0=16\pi^2(p-1)(q-1),\quad
 s=-4-v^2/2,\\
 \|\nabla^m\Rm\|_\infty\le C_m,\quad
 \|\nabla^m(\Ric+G_\ell)\|_\infty\le C_m/L.
 \end{gathered}
\]
常数依赖固定曲面、路径与阶数；它们不是所有夹逼度量的统一估计。
厚部极限为 $F\times\Sigma_{q-1,2}$，插值区体积 $O(\sqrt\ell)$，
注入半径趋零、直径发散；$K(e_0,e_3)=v^2/4$在非平凡剪切区为正。
因此构造不能套用要求处处负截面曲率的稳定性定理。
完整标架计算亦见\href{20260912.tex}{累积专论}第6章
\path{BUN:slow}、\path{BUN:ricci-matrix}、\path{BUN:slow-bound}。

同章还将慢插值扩展到经标准几何自由群商的单值表示
$\rho(a_i)=1,\rho(b_i)=[\phi_i]$，允许$\phi_i$不交换；
通过Moser面积规范和互不相交基底领圈保持同一构造机制。
该族签名仍为零。非零签名的任意曲面丛满足WP规范下
$\frac12\int|dF|^2\ge6\pi^2|\tau|$，故不能不改条件便把
固定曲率基底上的全部模参数速度一致压到零。

\subsection{高阶校正、TT障碍与退化逆算子}
以下为同一固定曲面丛上的进一步结论。证明见
\href{20260912.tex}{累积专论}第6章中所列命题；每个结论保留其背景与扰动尺度。
\begin{longtable}{L{.23\textwidth}L{.65\textwidth}}
\toprule 命题标签 & 结论及适用范围\\\midrule\endhead
\path{BUN:double} & 同时捏缩纤维$\varepsilon$和基底$\beta=\varepsilon^2$，采用归一化$\operatorname{sech}^3u$与$\operatorname{sech}t$轮廓，得$\delta=2\varepsilon^2/\pi^4+O(\varepsilon^4)$。端点Teichmüller距离也在退化，因此不违反固定端点的$L\delta$下界。\\
\path{BUN:third} & 求解三项张量校正的径向方程，核验Bianchi条件及尾项；截断后在同一闭流形上得每个固定$k$的$C^k$残差$O(\varepsilon^3)$。校正通常破坏原投影的黎曼下沉性。\\
\path{BUN:global-obstruction} & 真正全局TT测试张量给$c\varepsilon^3\le C\delta(G_\varepsilon+h)+C\varepsilon^2a+Ca^2$，$a=\|h\|_{C^2}$。$a=O(\varepsilon^2)$的校正具有三阶下界；这是指定邻域尺度障碍，不是任意度量的正缺陷间隙。\\
\path{BUN:conformal} & 对三阶背景$\widetilde G_\varepsilon$，整个共形类满足$\inf_{u>0}\delta(u^2\widetilde G_\varepsilon)=\Theta(\varepsilon^3)$。保留乘积开集还排除这些共形类中任意符号的Einstein度量；不排除其他共形类。\\
\path{BUN:volume-expansion} & 三阶背景的体积$V_0+\varepsilon^5/(2\pi^2)+O(\varepsilon^6\log(1/\varepsilon))$；负常标量归一化仍保留此首项。Yamabe共形类值与流形上确界不可混同。\\
\path{BUN:near-spectrum} & 原背景Jacobian负指数至少$6q-6$；校正后指定二维Ritz压缩的符号反转。两者是不同背景，有限压缩正定不证明整个算子可逆。\\
\path{BUN:inverse} & 在真正精确Einstein的闭乘积退化族上，去掉全部真实核后仍有$L^\infty$逆的对数增长，即使$L^2$正谱隙统一为正。这给出统一Newton推断的反向测试，不是Yau26反例。\\\bottomrule
\end{longtable}

\subsection{全局拓扑与变分不变量}
$E_D$为非球面流形，$\chi=4(p-1)(q-1)>0$、$\tau=0$、
$b_1=2p+2q-1$。其任何连通有限覆盖都无可积复结构且不承载非正截面曲率度量；
群论论证见\href{20260912.tex}{累积专论}的\path{BUNTOP:covers}。
这些不是一般Einstein不存在性定理。
局部Luttinger手术识别、覆盖单值、辛典范类、自由商保留操作以及
不能移入吹胀异常球的理由，完整保存于\path{BUNTOP:object}。

对$d$重连通覆盖，令$V_*=dV_0$。上述输入与monopole估计给
\begin{align}
 Y&=-4\sqrt{V_*},&
 \inf\int s^2&=16V_*,&\inf\int|\Ric|^2&=4V_*,\notag\\
 \inf_{s\ge-4}\Vol&=\inf_{\Ric\ge-g}\Vol=V_*,&
 \inf\int|\Rm_{\rm full}|^2&=8V_* .\label{eq:bundlefunctionals}
\end{align}
前述Yamabe/标量/Ricci/体积极值的等号迫使Kähler结构，故不取到；
最后一个全曲率泛函的等号恰等价于存在Einstein度量，达到性仍未知。
这不是矛盾：各泛函的等号条件不同。
两个取向全部monopole类的凸包和共同极化数值均可被同一输入序列饱和，
不能靠“再加入一些基本类”自动造出严格数值间隙。

假想精确Einstein度量必须为负常数，归一化后满足
\[
 V_*<\Vol(g)<3V_*,
\]
且必须有正截面曲率平面，两半Weyl在共同开稠密集上各有三个不同特征值。
命题\path{BUNTOP:Einstein-necessary}还给出连通归一化精确Einstein路径的体积恒定；
这只排除一路接近$V_*$的连通精确分支，不排除其他体积、其他分支或离散精确序列。

若\emph{另外假设}存在精确Einstein序列且体积趋$V_*$，
命题\path{BUNTOP:endpoint}、\path{BUNTOP:reducible-core}给出有界曲率、厚点、
非紧有限体积$\mathbb H^2\times\mathbb H^2$极限，排除不可约格点分支，
并给可约核心的整数周期、闭因子基本群单射及亏格整除条件。
候选中确有相容的乘积核心；未证明极限捕获全部体积或整个核心基本群单射。
这些是对\emph{假想精确序列}的条件退化结论，不能与已经存在的近似序列混写。

\subsection{下沉结构、谐和代表与源项}

适配某个纤维化、水平可积或基本平均曲率的Einstein刚性结果，可以限制\emph{适配度量}，却不能限制任意Einstein度量。曲面丛拓扑本身不要求未知Einstein度量尊重原纤维结构。

对基面NPC目标的能量极小调和映射，不自动是调和态射或水平共形。将两个双曲曲面的乘积投向另一个标记的双曲曲面，便可产生Einstein源上的调和映射而缺少所需共形性。优化目标Teichm\"uller结构给出的通常是积分TT平稳条件，不是逐点Hopf微分消失。欲获得全局障碍，应从具体无限Dehn-twist单值化、谱退化或共形缺陷找到新的强制量。

可以考察以下有数学裁决力的方向：是否某个任意Einstein度量都诱导特定可控的调和映射结构；是否可由单值化推出一个与Einstein方程不相容的积分恒等式；或是否另一种一般非Kähler构造实际上给$E_D$一个Einstein度量。对其他候选可以检查不同复结构，但$E_D$已经有上述无复结构结论，不能把这一已排除的方式重新包装成其开放可能。每一方向都需要完整假设，而不是“曲面丛没有复结构所以没有Einstein”。

与此相关的恒等式和反例见\href{20260912.tex}{累积专论}第7--10章：
\begin{itemize}
\item 第7章\path{WYL:source}：Weyl特征线的带源关系
$q=\tfrac32\ell_1-\eta H_1$及
$(9+4\eta^2)D-4\eta^2N=4|q|^2+8\eta\langle q,H_1\rangle$。
慢插值中源$q=O(L^{-1})$，而无源推断需要的相对量仅$O(L^{-2})$；
小散度不能当作精确调和Weyl。谐和压缩核也不能由近乎闭性删除。
\item 第8章\path{SUB:blocks}：带twist、非basic平均曲率和混合源项的完整下沉Ricci方程；
\path{SUB:integrable}、\path{SUB:basic}给水平可积或basic平均曲率时的乘积刚性。
\path{SUB:quartic}保留测度变分产生的四次矩与协方差，
\path{SUB:negative-covariance}给真实局部负协方差Einstein模型。
局部模型保持指定纤维几何的闭延拓另受命题\path{SUB:extension-gap}限制。
\item 第9章\path{SUB:harmonic-bochner}：任意假想$\Ric=-cG$下，谐和代表满足
$\int|\nabla df|^2+2\int J^2=c\int|df|^2$；等号才给乘积。
面积形式闭不等于余闭，Hodge投影不保持逐点可分解性，Euler配对看不到正合误差。
\item \path{SUB:optimal-target}：固定源存在唯一最优标记双曲目标，
但一阶条件只测试目标TT空间。
\path{SUB:optimal-nonconformal}给真正闭近Einstein族，其唯一最优目标映射仍不水平共形。
\item 第10章\path{GLB:WPbalance}：循环Dehn twist的WP长度方向使无强迫有效方程无解；
有源时却必须满足$\int w\,d\ell(R)=-\int w\operatorname{Hess}\ell(du,du)$。
没有证明完整Einstein方程给出足以禁止这项抵消的相对源控制。
\end{itemize}
这些反例区分了源项小、积分平稳、逐点共形与精确下沉之间不同的要求。

\section{全局变分问题与达到性}
四维有
\begin{equation}\label{eq:F}
\mathcal F(g)=\int_M(|W|^2+s^2/24)\dd
=8\pi^2\chi(M)+\frac12\int_M|E|^2\dd.
\end{equation}
如果某个\emph{固定}流形的输入构造还给出体积控制，使 $\int|E|^2\to0$，则其下确界为 $8\pi^2\chi$。若光滑度量达到这个下确界，则 $E=0$，由收缩Bianchi在 $n>2$ 得标量常数，从而是Einstein度量，符号不限。

但展示一列退化极小化序列，只证明一条极小化路径不紧，\emph{不证明不存在另一光滑极小值点}。近Einstein的 $L^p$ 残差也不是四阶变分Euler--Lagrange残差。Ekeland、Palais--Smale或“泛函非负”不能未经估计填平此差距。每次调用必须说明固定范畴、紧性、规范、允许退化和最终smooth type的恢复。

一个真正有影响的结果可以是：在这些已构造输入的流形上证明全局达到性，产生Einstein；或者证明所有可能达到情形都被一个严格障碍排除，得到全部Einstein不存在性。任一方向都要覆盖\emph{全部候选极小值}，而不是只覆盖选定orbifold附近的一支。

\section{正分支、覆盖与高维乘积}
$\Ric\approx+g$给Myers直径界、Bishop体积上界和有限基本群，
不提供统一体积下界或基本群阶上界。精确Einstein商本身也可以塌缩。

四维在固定非塌缩条件下可用有界Ricci正则片理论\cite{CN}；
几何有限型结论还需保持其完整假设。有限多个微分同胚型本身不构造Einstein度量，
也不使极限光滑，曲率集中与光滑结构恢复仍须处理。

若在一个非塌缩覆盖上产生了Einstein度量，仍要证明它对原甲板群等变，才能下降；覆盖存在Einstein并不自动意味着商的指定光滑结构存在所需度量。

对于 $n\ge5$，四维的 $\Lambda^2_\pm$、SW/Bauer--Furuta和CGB二次能量机制没有直接同构替代。高维路线可以寻找拓扑驱动的临界 $L^{n/2}$ 预算、特定特殊结构下的禁止曲率模式、沿余维四奇异层的ALE束障碍或新的全部Einstein不存在性。但若加入特殊holonomy、对称性、体积或完整曲率假设，只得到该范畴定理，不是原题。

高维统一机制的缺失阻止的是对所有维数的\emph{肯定证明}，
并非完整否定必须另过的一道高维关卡。
由式\eqref{eq:branches}，若在同一个固定维数 $n_0$ 同时否定
$P_{n_0}^+$与$P_{n_0}^-$，即可否定全维命题；两种符号可用不同流形序列。
单独的负分支反例仍不能否定 $Q_{n_0}$，而无需把两分支反例再传播到每个高维。

四维输入与同符号Einstein因子相乘，端同态误差是两因子误差的最大值，
因而可传播某些输入；它不传播全部Einstein不存在性，因为产品上的任意Einstein度量
不必分裂。正Einstein球因子维数至少2，这种乘积也跳过五维。


\section{不同端部模型的拓扑限制}
度二BCG间隙依赖双曲目标的Euler特征；由近等号熵映射取得更强结论，
还需要控制周边同伦信息与拓扑损失，不能直接把式\eqref{eq:bcg-threshold}变成维数统一界。

精确Poincar\'e--Einstein或AH/cusp模型显示切向与法向非线性补偿并非不可能，这是反对简单局部无解结论的真实依据。但一个自对偶模型未必提供可反射的全测地 $T^3$ 中心：若解析Einstein邻域中有反射，该反射逆转取向并交换 $W^+,W^-$，自对偶与反射相容会迫使二者同时消失；中心因而不能是所设想的紧致负常曲率三环面。这只限制该反射应用，不否定其他非对称中心。

已有负Einstein Dehn filling思路可以帮助研究边界与逆算子，但填充改变周边拓扑；得到另一个流形上的Einstein度量，不能成为Di Cerbo图流形上的存在反证。有限维的free-$T^2$、对角、orthogonally transitive或反射对称类也不穷尽所有图流形度量。保留这些路线的价值在于给主攻问题提供精确对照，而不是用局部范畴的失败宣布原题已经偏向真或假。

\section{开放的全局问题}
\subsection{双端连接器的连续极限}
固定完整Di Cerbo图流形候选、周边粘合和中心约束，构造保持两侧端点全部jets的完整矩阵场。数值上要求系数保存、精细网格独立验证、非线性去混叠、正定性下界、所有Ricci块和尺度依赖误差；再寻求解析或区间后验估计。若线性化逆范数为 $B$、二次非线性常数为 $C$、残差为 $\varepsilon$，校正可行性至少要控制相应 $B^2C\varepsilon$，而不是只让离散残差变小；余核与自由参数需另外解决。

若始终失败，寻找在分辨率和端点参数变化下仍存在的连续伴随障碍，而不是把一次优化停滞当成不存在性。有效障碍须对真实方程与允许自由度配对，并能排除所有足够小缺陷，而非只排除某个五次插值。

\subsection{曲面丛或固定解消上的一般Einstein度量}
对固定光滑类型，需要分别排除或构造正、零、负Einstein常数的度量。
可研究与该拓扑相容的调和代表、能量达到性或全局模空间；仅有限定背景附近的局部障碍
不足以决定全部Einstein不存在性。

如果证明了Einstein存在，应明确是新存在性结果而非反例失败的废弃产物。若只证明附近无解，保留为局部定理，并说明何种额外全局覆盖才能得到全部不存在性。

\subsection{图流形的Weyl能量与拓扑}
式\eqref{eq:graphW}使假想序列具有归一化Weyl小量。问题是这种全局近双曲状态
是否强迫相容的几何结构，或允许非双曲周边$\Z^3$在极小能量区域存活。
$L^2$小量不能直接升为全域点态Weyl小量。

这里的局部双曲模式确实存在；要取得全局障碍，还须控制周边群、粘合拓扑或
曲率集中区域的几何代价。

\section{旋转标架的有限区间障碍}
考虑一个比平切片更一般、但仍具有固定螺旋频率的度量类：
\[
 \omega_1=\cos(kz)dx+\sin(kz)dy,\qquad
 \omega_2=-\sin(kz)dx+\cos(kz)dy,
\]
\[
 g=dr^2+a(r)^2dz^2+b(r)^2\omega_1^2+c(r)^2\omega_2^2.
\]
$z$周期为$2\pi$，$k\in\Z$使余标架全局定义；$a,b,c>0$。
记
\[
 u=a'/a,\quad v=b'/b,\quad w=c'/c,\quad H=u+v+w,\quad
 d=\log(b/c),\quad p=d'=v-w,
\]
\[
 \mathcal K=2k^2a^{-2}\sinh^2d,\qquad
 \mathcal L=k^2a^{-2}\sinh(2d).
\]
在正交标架$(a\,dz,b\,\omega_1,c\,\omega_2)$中，
$d\omega_1=k\,dz\wedge\omega_2$、
$d\omega_2=-k\,dz\wedge\omega_1$给出切片截面曲率
\[
 K_{12}=-\mathcal K/2+\mathcal L,\qquad
 K_{13}=-\mathcal K/2-\mathcal L,\qquad K_{23}=\mathcal K/2.
\]
因而$\Ric_q=\operatorname{diag}(-\mathcal K,\mathcal L,-\mathcal L)$、
$R_q=-\mathcal K$。全Ricci为
\begin{align*}
 \Ric_{rr}&=-H'-(u^2+v^2+w^2),\\
 \Ric_{11}&=-\mathcal K-u'-Hu,&
 \Ric_{22}&=\mathcal L-v'-Hv,&
 \Ric_{33}&=-\mathcal L-w'-Hw.
\end{align*}
该不变标架中混合分量为零；这里没有省略Codazzi条件。
设$e_r,e_1,e_2,e_3$为$\Ric+3g$的四个对角分量，则
\begin{align}
 u'+Hu&=3-\mathcal K-e_1,&
 p'+Hp&=2\mathcal K\coth d-(e_2-e_3),\label{eq:helical-ode}\\
 (\log\mathcal K)'&=-2u+2p\coth d,&
 H^2-(u^2+v^2+w^2)&=6-\mathcal K+e_r-e_1-e_2-e_3.\nonumber
\end{align}

\begin{proposition}[反射中心后的不变区域]\label{prop:helical}
假设上述度量在$[0,R)$光滑，$r=0$为反射中心，
$u(0)=v(0)=w(0)=0$，且$\|\Ric+3g\|_{\op}\le\epsilon\le3/5$。
交换$b,c$使$d(0)>0$，则任意$r>0$有
\[
 u(r)<0,\qquad p(r)>0,\qquad
 \mathcal K(r)>\mathcal K(0)\ge6-4\epsilon.
\]
\end{proposition}
\begin{proof}
中心Hamilton约束给
$\mathcal K_0=6+e_r(0)-e_1(0)-e_2(0)-e_3(0)>0$，
所以$d(0)\ne0$。保留分量相消而不分别粗估，得到
$u'(0)=-3+e_2(0)+e_3(0)-e_r(0)\le-3+3\epsilon<0$。
在$p=0,d>0$的接触面上，
$p'>2\mathcal K-2\epsilon\ge12-10\epsilon>0$。
解由此进入$u<0,p>0$，而式\eqref{eq:helical-ode}使$\mathcal K$严格增加。
若首次重新达到$u=0$，便有
$u'=3-\mathcal K-e_1<3-\mathcal K_0+\epsilon\le-3+5\epsilon\le0$，
与首次离开矛盾；对$p=0$同理。端点$\epsilon=3/5$仍由
$\mathcal K>\mathcal K_0$获得严格号。证明不需要$H$的符号或完整性。
\end{proof}
因此任意正半径截面均有$\|S-I\|_{\op}>1$及
$\|\Ric_q\|_{\op}>\mathcal K\ge6-4\epsilon$；
不能达到形算子接近$I$、内禀曲率小的扩张端接口，即使退出半径随参数发散。
这只排除固定Bianchi VII$_0$旋转标架与反射中心的此一模型；
同时依赖两个底变量的轨道矩阵、非对角形算子或非反射中心不在该命题内。
计算可由\path{step78_bianchi_vii0_verify.py}逐项检验。

\section{正交传递连接器的定量障碍}
设完整四维度量具有自由、等距、正交传递的$T^2$作用，
商为$S^1\times\mathbb R$，写
\[
 g=\gamma_B+e^\sigma M_{ab}d\xi^ad\xi^b,
 \quad\det M=1,\quad
 U:B\to SL(2,\mathbb R)/SO(2),\quad e=|dU|^2.
\]
两端必须在紧集外\emph{精确向外扩张}：
$\gamma_B=dt_\pm^2+e^{2t_\pm}dz_\pm^2$、
$\rho=e^\sigma=c_\pm e^{2t_\pm}$、$U$常值。
假设$\|\Ric+3g\|_{\op,\infty}\le\epsilon$。
\begin{theorem}[完整双外扩端的OT障碍]\label{thm:OT}
在上述几何条件下，
\[
 12-8\epsilon\le
 \mathcal B(\epsilon):=
 \left[\frac{4\epsilon}{\sqrt{3+\epsilon}}
 +\sqrt{2(3+\epsilon)+\frac{20\epsilon^2}{3+\epsilon}}\right]^2.
\]
因此存在唯一正数$\epsilon_*$满足
$12-8\epsilon_*=\mathcal B(\epsilon_*)$，使$\epsilon<\epsilon_*$时不存在这种度量。
\end{theorem}
\begin{proof}
先构造极小分隔面。辅助度量$\widehat\gamma=\rho^2\gamma_B$在端部为
$c_\pm^2(e^{4t_\pm}dt_\pm^2+e^{6t_\pm}dz_\pm^2)$。
本原环路整体移向端部时长度发散；从固定紧集远行又返回的径向长度也发散。
故本原自由同伦类的极小化序列留在一个紧环带内，弧长参数化和下半连续性
给最短闭测地线；去除交叉段可取嵌入代表$C$。
由
$\Vol_3(\pi^{-1}C)=\operatorname{Vol}(T^2_{\mathrm{coord}})
 \int_C\rho\,ds_{\gamma_B}$的第一变分，其提升$\Sigma$极小。

取$\mathcal E=\Ric+3g$、$D=\tr\mathcal E-2\mathcal E(\nu,\nu)$、
$h=\partial_s\log\rho$。沿$C$记
$\langle f\rangle_C=(\int_C f\rho ds)/(\int_C\rho ds)$。
诱导标量曲率满足
$R_\Sigma=-2h_s-\frac32h^2-\frac12|U_s|^2$，而
$\int_C\rho h_sds=-\int_C\rho h^2ds$。
与Gauss方程$R_\Sigma=-6+D-|\mathrm{II}|^2$合并，得
\[
 \langle |U_s|^2\rangle_C=12-2\langle D\rangle_C
 +2\langle|\mathrm{II}|^2\rangle_C+\langle h^2\rangle_C
 \ge12-8\epsilon.
\]
故$\sup_Be\ge12-8\epsilon$。这里反射不是假设，两端向外扩张却不可省略。

再求能量上界。轨道方程给$|F|=|\tau_\sigma(U)|\le2\epsilon$。
取$L=\Delta+\langle\nabla\sigma,\nabla\cdot\rangle$，
真正的有限维带权Ricci为
\[
 \Ric^4_f=\Ric_B-\nabla^2\sigma-\tfrac12d\sigma^2
 =-3\gamma_B+\tfrac12U^*h_{\mathbb H^2}+\mathcal E_B
 \ge-(3+\epsilon)\gamma_B.
\]
置$A=U^*h_{\mathbb H^2}$、$V(X)=\langle dU(X),F\rangle$。
带权Bochner恒等式中目标曲率项为$2\det A$，且
$\langle dU,\nabla F\rangle=\operatorname{div}_\sigma V-|F|^2$。
由$\frac12\tr(A^2)+2\det A\ge\frac12e^2$得到
\[
 \tfrac12Le\ge |DdU|^2+\operatorname{div}_\sigma V
   -4\epsilon^2+\tfrac12e^2-(3+\epsilon)e.
\]
记$K=3+\epsilon$、$\kappa=2K$，配方$2\kappa e-e^2\le\kappa^2$给
\[
 (-L+\kappa)e\le\kappa^2+8\epsilon^2
                 -2\operatorname{div}_\sigma V.
\]
为处理散度源，令$p_t(x,y)$为$L$相对于$d\mu_\sigma=e^\sigma d\mu_B$的热核。
完整的$CD(-K,N)$空间上，$N=4$，热核质量为一，Fisher信息
$I(t)=\int p_t|\nabla_y\log p_t|^2d\mu_\sigma$满足
\[
 I'=-2\int p_t\Gamma_2(\log p_t)d\mu_\sigma
 \le2KI-\frac2N I^2.
\]
这里用到$\Gamma_2\ge-K\Gamma+(L\cdot)^2/N$、
$\int p_tL\log p_t=-I$及Jensen不等式。由正时间光滑化后比较Riccati方程，
\[
 I(t)\le\frac{NK}{1-e^{-2Kt}}.
\]
于是对$G_\kappa=\int_0^\infty e^{-\kappa t}p_t\,dt$，
Cauchy--Schwarz与$u=1-e^{-2Kt}$换元给
\[
 \sup_x\int|\nabla_yG_\kappa(x,y)|d\mu_\sigma(y)
 \le\frac{\sqrt{NK}}{2K}
       B\!\left(\frac{\kappa}{2K},\frac12\right)
 =\frac2{\sqrt{3+\epsilon}}.
\]
轨道能量和$V$在精确端部为零，故有界且散度分部积分无端部贡献。
对$(-L+\kappa)e$作热半群Duhamel比较，再令时间趋于无穷，
残余项$e^{-\kappa T}P_Te$消失。由于$|V|\le2\epsilon\sqrt e$，令$B_0=\sup e$，得到
\[
 B_0\le2(3+\epsilon)+\frac{4\epsilon^2}{3+\epsilon}
            +\frac{8\epsilon}{\sqrt{3+\epsilon}}\sqrt{B_0}.
\]
解关于$\sqrt{B_0}$的二次不等式得$B_0\le\mathcal B(\epsilon)$。
与分隔面下界合并即得结论。$\mathcal B$在非负半轴严格递增，
且$\mathcal B(0)=6<12$，故交点唯一并为正。
\end{proof}
数值$\epsilon_*\approx0.24625$可由\path{step87_resolvent_corrected.py}计算；
正间隙本身不依赖数值近似。在$\Ric\approx-g$规范下阈值除以3。
该定理不覆盖有联络twist、奇异轨道、一般非对称图流形度量或$S^3$连接颈。
有twist时，
twist张力导数的$-Qe$抵消漂移曲率中的$+Qe$，剩余项无定号。
极小分隔面下界本身仍可保留，失败的是未经重算的上界。

\section{带荷连接器：稳定性、局部构造与全局匹配}
与前一节相同，本节使用$\Ric\approx-3g$及
$\epsilon=\|g^{-1}\Ric+3\Id\|_{L^\infty,\op}$；
变成$\Ric\approx-g$规范时度量乘3、缺陷除3。

\subsection{共同变量：轨道形状、坐标荷与物理荷不同}
自由等距$T^2$作用的度量在主区域写成
\[
 g=\gamma_B+G_{ab}(d\xi^a+\mathcal A^a)(d\xi^b+\mathcal A^b),
 \quad G=\rho M,\quad\rho=e^\sigma,\quad\det M=1.
\]
$U:B\to\mathbb H^2=SL(2,\R)/SO(2)$描述$M$，目标度量取
$|dU|^2=\tfrac12\tr(M^{-1}dM\,M^{-1}dM)$。
若$d\mathcal A^a=f^a\,d\mu_{\gamma_B}$，定义
\begin{equation}\label{eq:charged-dictionary}
 \mathfrak q=\rho Gf=\rho^2Mf,\qquad
 Q=\tfrac12 f^{\mathsf T}Gf
   =\tfrac12 e^{-3\sigma}\mathfrak q^{\mathsf T}M^{-1}\mathfrak q,
 \quad L=\Delta+\langle\nabla\sigma,\nabla\cdot\rangle.
\end{equation}
$\mathfrak q$是坐标荷，$Q$是物理twist能量；前者常值不意味着后者常值。
记$\mathcal E=\Ric+3g$，避免与CGB中的无迹Ricci张量$E$混淆。
精确混合Einstein方程给$d\mathfrak q=0$，一般近似方程给
\begin{equation}\label{eq:physical-charge-source}
 e^{-3\sigma}M^{ab}\langle d\mathfrak q_a,d\mathfrak q_b\rangle
 =4|\mathcal E_{\mathrm{mix}}|_{\mathrm{HS}}^2,
 \quad
 \mathscr B=e^{-3\sigma/2}M^{-1/2}d\mathfrak q,
 \quad |\mathscr B|\le2\sqrt2\epsilon.
\end{equation}
这里保留完整物理范数；不能把它改成无权$d\mathfrak q=O(\epsilon)$。
精确Einstein连通区域若含一个精确OT开集，则常值荷为零。
仅有$Q\to0$的渐近端并不满足这条推论。

\subsection{Jacobi形式、容量与分隔面荷的高度界}
设$C$是$\rho^2\gamma_B$中最短本原闭环，提升为极小不变分隔面
$\Sigma=\pi^{-1}C$。沿$C$取$\gamma_B$弧长$s$、单位法向$\nu$，记
\[
 h=\sigma_s,\quad h_\nu=\partial_\nu\sigma,\quad
 e_s=|U_s|^2,\quad e_\nu=|U_\nu|^2,\quad
 V=|\mathrm{II}|^2=\kappa_C^2+\tfrac12h_\nu^2+\tfrac12e_\nu+Q,
 \quad D=\tr_g\mathcal E-2\mathcal E_{\nu\nu}.
\]
Gauss方程与内禀轨道计算分别给
\[
 R_\Sigma=-6+D-V,
 \qquad R_\Sigma=-2h_s-\tfrac32h^2-\tfrac12e_s,
\]
从而
\begin{equation}\label{eq:separator-h-identity}
 h_s=3-\tfrac12D+\tfrac12V-\tfrac34h^2-\tfrac14e_s.
\end{equation}
不能把极小性$H_\Sigma=0$替换成$\partial_\nu H_\Sigma=0$，
也不能由这条公式要求$R_\Sigma$为正。

最短环路的二阶变分对任意周期不变速度$\phi(s)$给
\begin{equation}\label{eq:separator-jacobi}
 \int_C\rho V\phi^2\,ds\le
 \int_C\rho\bigl[\phi_s^2+(3-\mathcal E_{\nu\nu})\phi^2\bigr]ds.
\end{equation}
一般情况下这只证明对不变变分稳定，不自动包括依赖环面变量的变分。
置$\phi=\rho^{-1/2}y$，展开导数并对交叉项分部积分，势为
$3-\mathcal E_{\nu\nu}-V+\frac12h_s+\frac14h^2$。
代入式\eqref{eq:separator-h-identity}得到无权形式
\begin{align}
 \int_C\mathcal W y^2\,ds
 &\le\int_C(y_s^2+c_{\mathcal E}y^2)\,ds,
 &\mathcal W&=\tfrac34V+\tfrac18(h^2+e_s),\label{eq:separator-schrodinger}\\
 c_{\mathcal E}&=\tfrac92-\mathcal E_{\nu\nu}-\tfrac14D,
 &\tfrac92-\tfrac32\epsilon&\le c_{\mathcal E}\le a_\epsilon:=\tfrac92+\tfrac32\epsilon.
\end{align}
取$y=1$并结合$\int_C h_s=0$还给
$\int_CV\,ds\le(24/5+8\epsilon/5)\operatorname{Length}(C)$。
这比仅有$\rho$加权平均多一项信息，但仍不是点态上界。

令$H_a=-\partial_s^2+a$为长度$\ell_C$的圆上的周期算子。
式\eqref{eq:separator-schrodinger}意味着
\[
 \|\mathcal W^{1/2}H_{a_\epsilon}^{-1}\mathcal W^{1/2}\|\le1.
\]
例如长度$\ell$区间的$H_a$容量为
\[
 \operatorname{Cap}_a(I_\ell)
 =a\ell+2\sqrt a\tanh\!\left(\frac{\sqrt a(\ell_C-\ell)}2\right).
\]
这里容量是$\int(y_s^2+ay^2)ds$在$y\ge1$于$I_\ell$上的下确界。
极小元在区间上为1，在补区间解$-y''+ay=0$且两端为1；
补区间长度$d=\ell_C-\ell$时，解为
$\cosh(\sqrt a(s-d/2))/\cosh(\sqrt a d/2)$。
将能量化成两端通量即得上述公式。
稳定二次型因此仍允许宽度趋零、质量有界而高度发散的势；
容量控制不能单独变成$\sup Q$控制。

物理Maxwell方程补充了变化率。令$w=\sqrt{2Q}$，在荷零点处以
范数的弱导数理解下式，避免除以$w$：
\[
 |w_s+\tfrac32hw|\le\tfrac12w\sqrt{e_s}+2\sqrt2\epsilon,
 \qquad
 |w_s|\le c(s)w+2\sqrt2\epsilon,
 \quad c^2\le\tfrac52(h^2+e_s)\le20\mathcal W.
\]
其直接推导是对$w^2=e^{-3\sigma}\mathfrak q^{\mathsf T}M^{-1}\mathfrak q$求导。
矩阵$M^{-1/2}M_sM^{-1/2}$无迹，且算子范数等于$\sqrt{e_s}$；
形状导数贡献至多$\frac12w\sqrt{e_s}$，荷导数贡献至多$|\mathscr B_s|$。
先在$w>0$处计算，再以$\sqrt{w^2+\eta^2}$逼近即可得到零点处的弱不等式。
当$\ell_C\ge2/\sqrt{a_\epsilon}$，以荷最大点为中心取
$r=a_\epsilon^{-1/2}$的帐篷试验函数，得到半径$r/2$弧上
$\int\mathcal W\le(32/3)\sqrt{a_\epsilon}$。
该弧存在$Q\le128a_\epsilon/9$的一点，连接到最大点时
$\int c\le\sqrt{320/3}$；Gronwall遂给
\begin{equation}\label{eq:separator-Q-bound}
 \sup_CQ\le C_{\mathrm{long}}(\epsilon):=
 \frac12e^{2\sqrt{320/3}}
 \left(\frac{16}{3}\sqrt{a_\epsilon}
       +\frac{\sqrt2\epsilon}{\sqrt{a_\epsilon}}\right)^2.
\end{equation}
短圆情形用$\int\mathcal W\le a_\epsilon\ell_C$，从一个
$Q\le4a_\epsilon/3$的点沿较短弧积分即可；同一计算给可用粗界
\[
 C_{\mathrm{short}}(\epsilon)=\tfrac12e^{4\sqrt{10}}
 \left(\sqrt{8a_\epsilon/3}
       +\frac{2\sqrt2\epsilon}{\sqrt{a_\epsilon}}\right)^2.
\]
因此$\sup_CQ\le\max\{C_{\mathrm{long}},C_{\mathrm{short}}\}$不依赖圆长
或$\rho$振荡。它很大，且$\epsilon\to0$时不趋零；它也没有给
$\sup_BQ$；后者还需要法向传输估计。

\subsection{全稳定的精确Koiso芽与法向能量集中}
这里所用的解析Cauchy输入可以明确表述：解析三维度量$q$与解析对称算子
$\mathsf S$若满足
\[
 \operatorname{div}_q\mathsf S-d\tr\mathsf S=0,\qquad
 R_q+\tr\mathsf S^2-(\tr\mathsf S)^2=2\lambda,
\]
则在$\Sigma\times\{0\}$附近存在唯一解析Einstein芽
$g=dn^2+q_n$，$\Ric_g=\lambda g$，给定$q_0=q$及
$\dot q_0=-2q(\mathsf S\cdot,\cdot)$。原定理的Weingarten约定为
$\nabla_X\nu=-\mathsf S X$，故与前文$K=\dot q/2$差一个整体符号。
其证明以切向Einstein方程的解析Cauchy--Kowalevskaya解开始，
由Bianchi关系传播约束，再由唯一性拼合局部片；紧$\Sigma$允许为一个固定初值
选共同邻域，但不提供对一列退化初值统一的厚度。光滑而非解析的数据不能直接套用。
这就是Ammann--Moroianu--Moroianu所述Koiso定理\cite[Theorem 2.1]{Ammann}；
相关输入也见
\href{20260912.tex}{\texttt{20260912.tex}}的\path{FND:analyticcollars}。
在圆上选正解析函数$P(s)$，取闭解析浸入
$U:S^1\to\mathbb H^2$使$|U_s|^2=12+2P$；
具体令$L_U=\int_0^{\ell_C}\sqrt{12+2P(s)}ds$，
选择双曲圆半径$R_U$使$2\pi\sinh R_U=L_U$，
再取角速度$\theta_s=\sqrt{12+2P}/\sinh R_U$。
它一周恰转$2\pi$，给全局解析目标曲线及其单位法向。
令$N_U$为目标曲线的单位法向，置
\[
 q_\Sigma=ds^2+M(U(s))_{ab}dx^adx^b,
 \quad A=\tfrac12M^{-1}M_s,
 \quad B=\sqrt{P/2}\,M^{-1}dM(N_U),
 \quad \mathsf S=\operatorname{diag}(0,B).
\]
这里$\nabla_X\nu=-\mathsf S X$；换法向将$\mathsf S$整体改号。
目标度量规范给$\tr B=0$、$\tr B^2=P$、$\tr(AB)=0$，故
\[
 R_{q_\Sigma}=-6-P,\qquad\tr\mathsf S=0,
 \quad\tr\mathsf S^2=P,\qquad
 \operatorname{div}\mathsf S-d\tr\mathsf S=0.
\]
这些是$\Ric=-3g$的完整标量和动量Cauchy约束；对应
$U_\nu=\pm\sqrt{2P}\,N_U$，所以$e_\nu=2P$。

取解析尖峰$P_k\,ds\rightharpoonup m\delta_0$，其中
$m<2\sqrt3\tanh(\sqrt3\ell_C/2)$。
例如可取
\[
 P_k(s)=\frac{m\,e^{k\cos(2\pi s/\ell_C)}}
                 {\int_0^{\ell_C}e^{k\cos(2\pi t/\ell_C)}dt}.
\]
$-\partial_s^2+3$的周期Green核在对角线上等于
$\coth(\sqrt3\ell_C/2)/(2\sqrt3)$。
由评价泛函的Riesz表示，
\[
 |y(0)|^2\le
 \frac{\coth(\sqrt3\ell_C/2)}{2\sqrt3}
 \int_0^{\ell_C}(y_s^2+3y^2)ds.
\]
严格质量条件给点势的正间隙；$H^1(S^1)$有界集一致Hölder连续，
故$\int P_ky^2\to m|y(0)|^2$在单位能量球上一致成立。因此充分大$k$时
$-\partial_s^2+3-P_k>0$。
本构造在环面方向平移不变，非零Fourier模增加非负势
$\mathbf k^{\mathsf T}M^{-1}\mathbf k$，故这里不只是“不变稳定”，
而是对全部变分严格稳定。
由解析Einstein Cauchy存在定理，每个固定$k$的解析约束数据
产生一个真正的局部Einstein芽，同时$\sup e_\nu=2\sup P_k\to\infty$。
环面平移保持初值，解析唯一性使该作用延拓为芽上的等距作用。
其存在邻域可随$k$缩小，不能据此得到完备度量、共同厚度的领圈或双端连接器。

要将分隔面荷界向法向传输，Maxwell方程仍需控制
\[
 \int_\Gamma\bigl(|\partial_\nu\sigma|+\sqrt{e_\nu}\bigr)\,dn.
\]
单层Jacobi容量与轨道方程的二阶迹约束不提供该量的共同上界。
这一构造保留了一个具体障碍：即使局部精确Einstein、完整Cauchy约束和全部变分稳定
同时成立，仍不足以由单层数据取得统一的法向能量控制。

\subsection{非零荷的稳定中心与二变量局部展开}
局部约束也不强迫twist消失。固定$0<S<3$，取解析目标圆
$|U_s|^2=12+2S$，选择非零常值荷$\mathfrak q$足够小，使
$Q(s)=\frac12\mathfrak q^{\mathsf T}M^{-1}\mathfrak q<S$。
记$P=S-Q>0$，令
\[
 B=\sqrt{P/2}\,M^{-1}dM(N_U),\quad m=-\tfrac12\mathfrak q,
 \qquad
 \mathsf S=\begin{pmatrix}0&m_b\\m^a&B^a{}_b\end{pmatrix},
 \quad m^a=M^{ab}m_b.
\]
混合块在范数中出现两次，因此
\[
 \tr\mathsf S=0,\qquad\tr\mathsf S^2=P+2m^{\mathsf T}M^{-1}m=P+Q=S,
 \quad R_{q_\Sigma}=-6-S.
\]
动量方程的水平分量为$-\tr(AB)=0$；垂直分量中
\[
 m_a'+(\tr A)m_a-A^b{}_am_b+(MA)_{ba}m^b=0,
\]
因为$m$常值、$\det M=1$，末两连接项相消。
法向或曲率约定改变时荷与$m$的符号须一并改变，平方量和约束不变。
中心的完整Jacobi形式为
\[
 \mathcal J(\phi)=\int_\Sigma
       \bigl(|\nabla\phi|^2+(3-S)\phi^2\bigr)d\mu,
\]
故对全部变分严格稳定，最低本征值为$3-S$。
解析Cauchy定理给具有非零保守荷的局部精确Einstein芽；它不能在同一连通精确区域
直接接到荷为零的精确OT开集，后者需要混合Ricci源或改变问题的端部条件。

为保留实际局部研究而非只留一个中心，写
$g=dn^2+q(n,s)$、$K=\frac12q_n$。切向方程为
\[
 q_n=2K,\qquad K_n=\Ric(q)+3q-HK+2Kq^{-1}K.
\]
每层分解
\[
 q=a^2ds^2+G_{ab}(dx^a+\mathcal A_s^a ds)(dx^b+\mathcal A_s^b ds).
\]
若$v_a=q_{as}$，则
$\mathcal A_s=G^{-1}v$、$a^2=q_{ss}-v^{\mathsf T}G^{-1}v$、
$\rho=\sqrt{\det G}$。在$\mathcal A_n=0$规范，
$f=(\partial_n\mathcal A_s)/a$，由式\eqref{eq:charged-dictionary}重建荷。
在上述中心$H=0,\rho=a=1,|q^{-1}K|^2=S$处，正常方程与Schur补给
\begin{equation}\label{eq:charged-taylor}
 H_n=3-S,\qquad a_n=0,\quad a_{nn}=-3-S-Q(s),
 \qquad\sigma_n=0,\quad\sigma_{nn}=6+Q(s).
\end{equation}
所以底圆先收缩而轨道面积先扩张；两侧按各自向外法向都有初始体积扩张。
这不是所有系数只依赖$n$的平切片模型，后者仍受$\epsilon\ge3/2$障碍。

完整$3\times3$对称张量与周期Fourier离散可用于检验Taylor展开。
例如$S=1,n=0.02$时，分辨率$48,64,96$的标量／动量残差分别为
$(8.32\cdot10^{-2},5.79\cdot10^{-2})$、
$(9.96\cdot10^{-3},1.01\cdot10^{-2})$、
$(1.32\cdot10^{-4},2.29\cdot10^{-4})$。
这组计算由\path{step96_koiso_einstein_cauchy_evolution.py}实现，不给出固定正长度延拓。
Einstein方程按法向当作演化仍是椭圆Cauchy问题，长距离高频增长不能用加步数消除。

\subsection{两荷恒等式、Noether配对与红移反模型}
固定非零荷时，令
$F(U)=\frac12\mathfrak q^{\mathsf T}M^{-1}\mathfrak q$、$\beta=\log F$。
在上述双曲目标规范，直接微分给
\[
 \operatorname{Hess}F=Fh_{\mathbb H^2},\quad|\nabla F|^2=F^2,
 \quad\operatorname{Hess}\beta=h_{\mathbb H^2}-d\beta\otimes d\beta,
 \quad|\nabla\beta|=1.
\]
为核对这些等式，可作常值$SL(2,\R)$基变换使
$\mathfrak q=(c,0)$。在
$M=\left(\begin{smallmatrix}
e^\alpha+\zeta^2e^{-\alpha}&\zeta e^{-\alpha}\\
\zeta e^{-\alpha}&e^{-\alpha}
\end{smallmatrix}\right)$、
$h_{\mathbb H^2}=d\alpha^2+e^{-2\alpha}d\zeta^2$中，
$F=(c^2/2)e^{-\alpha}$；计算两个对角Hessian分量及混合分量即可。
精确面积与轨道方程为$L\sigma=6+Q$、
$\tau_\sigma(U)=-Q\nabla\beta$。由$\log Q=-3\sigma+\beta\circ U$得到
\begin{equation}\label{eq:charged-logQ}
 L\log Q=-18-4Q+e_\perp,
 \qquad e_\perp=|dU|^2-|d(\beta\circ U)|^2\ge0.
\end{equation}
仅荷对齐的对角模型才有$e_\perp=0$。
即使如此，在正最大值处$LQ=-18Q-4Q^2<0$也符合最大值原理，绝非矛盾。
一维方程$y''=-18-4e^y$具有
$(y')^2=36(\log Q_0-y)+8(Q_0-e^y)$，
允许$Q=e^y$在中央取正最大值而两端趋零。
这是符号反模型，不是完整Einstein连接器。

变量荷下，残差块定义为
\begin{align}
 \mathscr H&=\Ric_B-\operatorname{Hess}\sigma
       -\tfrac12d\sigma^2-\tfrac12U^*h_{\mathbb H^2}-Q\gamma_B+3\gamma_B,\nonumber\\
 \mathscr T&=L\sigma-6-Q,\qquad
 \mathscr Z=\tau_\sigma(U)+\nabla_UQ,
 \qquad m_{\mathfrak q}=e^{-3\sigma}
          \langle M^{-1}\mathfrak q,d\mathfrak q\rangle.\label{eq:charged-blocks}
\end{align}
利用收缩Bianchi恒等式与调和映射应力散度，得到
\begin{equation}\label{eq:charged-noether}
 2\operatorname{div}_\sigma\mathscr H-d\tr\mathscr H
 =-d\mathscr T-\mathscr T\,d\sigma
 -\langle\mathscr Z,dU\rangle-m_{\mathfrak q}.
\end{equation}
具体计算为
\[
 2\operatorname{div}_\sigma\mathscr H-d\tr\mathscr H
 =-d(L\sigma)-(L\sigma+2Q-6)d\sigma
                  -\langle\tau_\sigma(U),dU\rangle.
\]
其中$2\operatorname{div}\Ric-dR=0$，
$\operatorname{div}(U^*h)=\langle\tau(U),dU\rangle+\tfrac12d|dU|^2$，
且对Hessian和$d\sigma^2$求散度后合成$d(L\sigma)$。
最后代入
$L\sigma=\mathscr T+6+Q$、
$\tau_\sigma(U)=\mathscr Z-\nabla_UQ$及
$dQ=-3Qd\sigma+\langle\nabla_UQ,dU\rangle+m_{\mathfrak q}$，
三项$Qd\sigma$与目标梯度项相消，留下式\eqref{eq:charged-noether}。
对紧支撑$X$和$d\mu_\sigma=\rho d\mu_B$分部积分，得到
\begin{align}
 \int_Bm_{\mathfrak q}(X)d\mu_\sigma
 =\int_B\bigl[&2\mathscr H:\operatorname{sym}\nabla X
  -(\tr\mathscr H)\operatorname{div}_\sigma X\nonumber\\
 &+\mathscr T\operatorname{div}X
  -\langle\mathscr Z,dU(X)\rangle\bigr]d\mu_\sigma.\label{eq:charged-pairing}
\end{align}
它把导数从$C^0$小的残差移到测试场，却不自动控制测试场、轨道能量或加权体积。
$\mathscr H,\mathscr T,\mathscr Z,d\mathfrak q$属于同一个Ricci残差，不能独立指定。

径向专门化$\gamma=dr^2+\psi^2d\theta^2$、$\mathscr H=C\,dr^2$、
$u=\sigma',j=\psi'/\psi$，当$\mathscr T=\mathscr Z=0$时，式\eqref{eq:charged-noether}给
\[
 C'=-2(u+j)C-wb,
 \qquad\int e^{2\sigma}\psi^2wb\,dr=0
\]
（后式还需要两端$C=0$）。这个第二加权矩不同于坐标荷的
$\int\mathfrak q'\,dr=0$，说明只消去总荷通量不足以完成水平匹配。

\paragraph{面积方程也不能单独阻止红移。}
考虑如下标量子系统：令$q=1$在$|r|\le R$，
再以固定宽度、端点平坦的光滑截断降到$q=0$；取
\[
 \rho''=(6+Q)\rho,\quad Q=\tfrac12\rho^{-3}q^2,
 \quad\rho(0)=1,\ \rho'(0)=0,
 \quad b=\rho^{-3/2}q',\quad w=\rho^{-3/2}q.
\]
于是$w'+\frac32\sigma'w=b$、两端$q=0$与面积方程同时成立。
中央平台的一阶积分为
\[
 (\rho')^2=6\rho^2-\rho^{-1}-5,
\]
故$\rho\asymp e^{\sqrt6r}$，将截断推远使
\[
 \|b\|_\infty\asymp e^{-3\sqrt6R/2},
 \qquad
 \frac{\int w^2\rho\,dr}{\int b^2\rho\,dr}\asymp e^{2\sqrt6R},
 \qquad Q(0)=\tfrac12.
\]
这否定仅由Maxwell、两端零荷、面积方程与分隔面有限荷界推出无迹项的统一强制估计。
它不是几何连接器：若冻结$M=I$，$Q>0$时轨道方程仍有未满足的源，
水平方程也未闭合。任何有效的排除或构造，必须使用这些尚未满足的耦合方程。

\subsection{局部恒等式与共享参数的双端匹配算子}
式\eqref{eq:charge}中的矩阵此处为$A=e^{-3\sigma}M^{-1}$。
对固定荷精确芽和紧支撑$X$，
\[
 \tfrac12\mathfrak q^{\mathsf T}
  [X(A)+(\operatorname{div}_\sigma X)A]\mathfrak q
 =\operatorname{div}_\sigma(QX)-m_{\mathfrak q}(X).
\]
因$d\mathfrak q=0$，积分恒为零，无须使用Einstein方程。
所以局部$K_X$数值接近零只检查分部积分和荷保存，不能作为独立的全局阻碍，
也不是又一次Noether方程验证。可能有效的测试场必须与两端及同一中央数据关联。

具体全局对象是商柱面$\Omega=\R_\tau\times S^1_s$上的混合阶问题。
取$\gamma=e^{2\varphi}(d\tau^2+ds^2)$，完整保留$s$模态；未知量为
$(\varphi,\sigma,U,\mathfrak q,\mathcal A;p)$，其中$p$包含
\emph{两侧共用}的中心$S$、荷向量、目标相位／标架，以及两端尺度和轨道形状。
$\mathcal A$须固定规范并保留两端平坦联络holonomy，不能只保留局部曲率。
残差映射采用
\[
 \mathscr E=(\mathscr H,\mathscr T,\mathscr Z,\mathscr B)
\]
及式\eqref{eq:physical-charge-source}、\eqref{eq:charged-blocks}的物理规范。
荷是\emph{一阶}变量，不能虚构为另一个二阶sigma-model场。
共形／Coulomb规范下，其主部为
\begin{align}
 D\mathscr T_2&=e^{-2\varphi}\Delta_0\dot\sigma,&
 D\mathscr Z_2&=e^{-2\varphi}\Delta_0\dot U,\nonumber\\
 D(\tr\mathscr H)_2&=e^{-2\varphi}
       (-2\Delta_0\dot\varphi-\Delta_0\dot\sigma),&
 D\mathscr B_1&=e^{-3\sigma/2}M^{-1/2}d\dot{\mathfrak q}+\cdots.
\end{align}
压去非零共形权重后，对非零余向量$\xi=(\xi_1,\xi_2)$，
令$r^2=\xi_1^2+\xi_2^2$，该混合阶主符号为
\[
 P(\xi)=
 \begin{pmatrix}
 -2r^2&-r^2&0&0&0&0\\
 0&r^2&0&0&0&0\\
 0&0&r^2&0&0&0\\
 0&0&0&r^2&0&0\\
 0&0&0&0&i\xi_1&0\\
 0&0&0&0&i\xi_2&0\\
 0&0&0&0&0&i\xi_1\\
 0&0&0&0&0&i\xi_2
 \end{pmatrix}.
\]
前四列对应$(\varphi,\sigma,U_1,U_2)$，后两列对应荷。
于是方形二阶子块行列式为$-2|\xi|^8$，
而$\det(P^*P)=4|\xi|^{20}>0$。
这证明所列主部单射，不证明全局Fredholm性质、边界补条件、指数、
实际余核或一致反算子；Bianchi冗余和联络重建仍须在完整问题中正确处理。

一种两半Dirichlet--Neumann表述是：分别从两端求解到中央领圈，要求返回数据
同时等于同一个$\mathsf S(S,\mathfrak q,\text{相位})$。
有限匹配对象至少包括两分量荷的生成与消去、两端holonomy、
中央形状和水平Dirichlet／Neumann条件，以及完整算子的真正伴随模投影。
可用的投影候选包括水平无迹张量、与两端关联的Noether场、目标$\mathbb H^2$的
三个Killing流；它们不是预先独立可调的参数，也不是已算出的余核基。

若一个固定荷桥确实延拓到向外扩张区，远处长度$\ell$的截断可能具有
$\|\mathscr B\|_\infty=O(e^{-3R}/\ell)$、$\|Q\|_\infty=O(e^{-6R})$。
这只是条件性红移方向，尚未构造那座桥。
完整Fourier--Chebyshev线性化的最小奇异值、对应左模及其随$R$的去向，
可区分逃向无穷远的近核与中央耦合的真实有限障碍；数值诊断仍须回到连续算子的误差控制。
上述局部构造、稳定性和全局匹配框架均未解决所有free-$T^2$近似连接器，
更没有排除非对称度量或完成原题两符号分支。

\section{计算实现}
相关程序全文收录于\href{../references/computations.md}{计算程序附录}，
相应数值输出见\href{../references/experiments.md}{计算结果附录}；
两者均可按下列文件基名检索。符号程序用于展开和核对正文恒等式，
离散演化用于考察解析方程的局部行为；几何结论的假设和推导由正文给出。
\begin{longtable}{L{.57\textwidth}L{.31\textwidth}}
\toprule 文件基名 & 数学计算\\\midrule\endhead
\path{step91_sambusetti_bcg_pinching_threshold} & 度二映射的熵--体积约束与Ricci夹逼阈值。\\
\path{step78_bianchi_vii0_verify} & 旋转标架的切片曲率与Einstein约化方程。\\
\path{step87_resolvent_corrected} & OT模型的resolvent常数与正阈值。\\
\path{step92_weighted_jacobi_transform} & 共轭后的各项系数。\\
\path{step92_jacobi_maxwell_bound} & Jacobi容量与Maxwell荷的长圆估计。\\
\path{step92_constrained_spike_koiso} & 尖峰中心、动量约束及规范。\\
\path{step94_invariant_charge_identity} & 双曲目标Hessian与荷的链式法则。\\
\path{step94_two_ended_redshift_taper} & 标量子系统、衰减率及未满足方程。\\
\path{step94_charged_koiso_separator} & 带荷中心和稳定性系数。\\
\path{step94_horizontal_noether_identity} & 水平Bianchi恒等式的系数检查。\\
\path{step96_initial_taylor_invariants} & 中心Taylor、Schur补与荷重建。\\
\path{step96_relative_KX_diagnostic} & 固定荷局部空检验。\\
\path{step96_linear_matching_symbol} & 实现范围内的混合阶主符号。\\\bottomrule
\end{longtable}
解析Cauchy方程的短区间离散演化使用
\path{step96_koiso_einstein_cauchy_evolution.py}；
细网格输出为
\path{step96_koiso_einstein_cauchy_evolution_N96_short.json}，
分别在上述两份附录中检索。其中央Taylor系数、约束和变量重建公式已在正文展开，
离散演化不承担双端存在性证明。

\begin{thebibliography}{99}
\bibitem{Yau1982} S.-T. Yau, \emph{Problem Section}, in \emph{Seminar on Differential Geometry}, Annals of Mathematics Studies 102, Princeton University Press, 1982, pp. 669--706.

\bibitem{BCG} G. Besson, G. Courtois and S. Gallot,
\emph{Entropies et rigidités des espaces localement symétriques de courbure strictement négative},
Geom. Funct. Anal. 5 (1995), 731--799.

\bibitem{Ammann} B. Ammann, A. Moroianu and S. Moroianu,
\emph{The Cauchy problems for Einstein metrics and parallel spinors},
arXiv:1106.2066. \url{https://arxiv.org/abs/1106.2066}.

\bibitem{DiCerbo} L. F. Di Cerbo, \emph{Extended graph 4-manifolds, and Einstein metrics}, arXiv:2106.13279. \url{https://arxiv.org/abs/2106.13279}.

\bibitem{RT24} J. G. Ratcliffe and S. T. Tschantz,
\emph{Hyperbolic 24-cell 4-manifolds with one cusp},
arXiv:2004.07284, \S\S9,11 and Table 2.
\url{https://arxiv.org/abs/2004.07284}.

\bibitem{CN} J. Cheeger and A. Naber, \emph{Regularity of Einstein manifolds and the codimension 4 conjecture}, arXiv:1406.6534. \url{https://arxiv.org/abs/1406.6534}.

\bibitem{NZ} A. Naber and R. Zhang, \emph{Topology and $\epsilon$-regularity theorems on collapsed manifolds with Ricci curvature bounds}, Geom. Topol. 20 (2016), 2575--2664. \url{https://doi.org/10.2140/gt.2016.20.2575}.

\bibitem{JKX} W. Jiang, L. Kong and S. Xu, \emph{Convergence of Ricci-limit spaces under bounded Ricci curvature and local covering geometry I}, arXiv:2205.00373. \url{https://arxiv.org/abs/2205.00373}.

\bibitem{Anderson1989} M. T. Anderson, \emph{Ricci curvature bounds and Einstein metrics on compact manifolds}, J. Amer. Math. Soc. 2 (1989), 455--490. \url{https://doi.org/10.1090/S0894-0347-1989-0999661-1}.

\bibitem{GJ2017} H. Ge and W. Jiang, \emph{$\epsilon$-regularity for shrinking Ricci solitons and Ricci flows}, arXiv:1707.05511. \url{https://arxiv.org/abs/1707.05511}.

\bibitem{OzuchII} T. Ozuch, \emph{Noncollapsed degeneration of Einstein 4-manifolds, II}, arXiv:1909.12960v2; Geom. Topol. 26 (2022). \url{https://arxiv.org/abs/1909.12960}.

\bibitem{LeBrun2001} C. LeBrun, \emph{Ricci curvature, minimal volumes, and Seiberg--Witten theory}, Invent. Math. 145 (2001), 279--316. \url{https://arxiv.org/abs/math/0003068}.

\end{thebibliography}
\end{document}
